% GENERATED FROM typst/manual/tables/swells_k1_k2_k3_comparison_table.typ BY scripts/typ2tex.py (v1.0.0) — DO NOT EDIT.
% Edit canonical Typst under typst/ and regenerate with: python3 scripts/typ2tex.py --root . --outdir latex

\begin{table}[H]
\centering
\caption{SWELLS K1/K2/K3 diagnostic comparison.}
\label{tab:manual-lensing-swells-k123}
{\footnotesize
\begin{tabularx}{\linewidth}{l>{\raggedright\arraybackslash}p{0.280\linewidth}cccl}
\toprule
Run & Input basis & $\theta_E$ $\delta_3$ $m_{\mathrm{tot}}$ & $\theta_E$ match [\%] & $\sigma$ $\delta_3$ $m_{\mathrm{tot}}$ & Decision \\
\midrule
K1 & $\sigma_{\mathrm{SDSS}}$ + H$\alpha$ + axis ratio & 8 & 40.0 & 4 & retained baseline \\
K2 & K1 + redshift geometry context & 18 & 40.0 & 4 & audit diagnostic \\
K3 & Paper I Table 2 bulge/disk medium fields & 10 & 33.3333 & 12 & audit diagnostic \\
\bottomrule
\end{tabularx}
}
\end{table}

\par\vspace{0.4em}\par

\emph{K1 reason.}Lowest retained $\delta_3$ residual. 

\emph{K2 reason.}Over-amplified high-$\theta_E$ class. 

\emph{K3 reason.}Under-resolved observable map with current medium data.
